%% kapa: exam=1 round=26 year=2015 period=1 form=A subject=economics subject_ko=경제원론 seq=19 number=59
%% source: Q-Net 감정평가사 기출문제 (cst00309) artlSeq=5138001 file=extracted/1008/5138001/제26회_감정평가사 1차_기출/제26회_감정평가사 1차_기출/1교시_감평_A형.pdf
%% conversion: kordoc 4.24.51 -> tools/gen1.py (auto); review=proofread-sonnet; flags=none
\begin{question}{59}
다음 표는 이동통신시장을 양분하고 있는 甲과 乙의 전략(저가요금제와 고가요금제)에 따른 보수행렬이다. 甲과 乙이 전략을 동시에 선택하는 일회성 게임에 관한 설명으로 옳지 않은 것은? (단, 괄호 속의 왼쪽은 甲의 보수, 오른쪽은 乙의 보수를 나타낸다.)
\par{\small\centering\begin{tabular}{c|c|c|c|}
\multicolumn{2}{c}{} & \multicolumn{2}{c}{乙} \\ \cline{3-4}
\multicolumn{2}{c|}{} & 저가요금제 & 고가요금제 \\ \cline{2-4}
\multirow{2}{*}{甲} & 저가요금제 & (500, 500) & (900, 400) \\ \cline{2-4}
 & 고가요금제 & (300, 800) & (700, 600) \\ \cline{2-4}
\end{tabular}\par}
\begin{choices}
  \item 甲은 乙의 전략과 무관하게 저가요금제를 선택하는 것이 합리적이다.
  \item 甲이 고가요금제를 선택할 것으로 乙이 예상하는 경우 乙은 고가요금제를 선택하는 것이 합리적이다.
  \item 甲과 乙의 합리적 선택에 따른 결과는 파레토 효율적이지 않다.
  \item 내쉬균형(Nash equilibrium)이 한 개 존재한다.
  \item 乙에게는 우월전략이 존재한다.
\end{choices}
\end{question}
\answer{2}
